%% =================================================================
%%  1.  Introduction
%% =================================================================
\section{Introduction}
\label{sec:intro}

A quasiperfect number (QPN) is a positive integer~$N$ whose sum of
divisors satisfies $\sigma(N) = 2N + 1$.  While no such numbers have
been found, their existence has been heavily constrained by decades of
mathematical research.  Cattaneo (1951) demonstrated that any QPN must
be an odd perfect square~\cite{cattaneo1951}.  Subsequent computational
and analytical efforts by Hagis and Cohen (1982) established that a QPN
must be greater than $10^{\TelemetryMinLog}$ and contain at least seven distinct prime
factors~\cite{hagis1982}.  Despite these bounds, closing the search
space remains computationally difficult due to combinatorial explosions
in possible prime factorizations.

\paragraph{The verification gap.}
The central obstacle confronting modern computational attacks on the QPN
problem is not raw arithmetic speed: it is the \emph{verification gap}.
Purely mathematical approaches struggle to eliminate vast swaths of
candidate numbers without concrete enumeration.  Purely computational
searches, meanwhile, are prone to unverified optimisations and integer
overflow bugs that might silently skip a valid QPN\@.  Translating
formally proved bounds into executable code introduces a fresh class of
risks---floating-point inaccuracies, off-by-one errors in sieve
implementations, and race conditions in parallel search---that can
compromise mathematical soundness without any visible failure.

\paragraph{Our contribution.}
The Unified Algebraic-Lattice Bipartition Framework (UALBF) attacks the
verification gap head-on.  Instead of presenting mathematics and
computation as separate concerns, the UALBF \emph{fuses} them: the Rust
search engine calls compiled Lean~4 binaries through a C-level Foreign
Function Interface (FFI)\@.  Critical hot-path operations---128-bit
modular inverses for ray-cast targeting and exact $\sigma(p^{2e})$
evaluations---are dispatched into Lean code whose correctness is
established by formal bridge theorems: \texttt{computeSigmaNat\_eq\_sigma}
proves that the FFI's $\sigma$ computation equals Mathlib's
sum-of-divisors function, and \texttt{extGcdAux\_bezout} establishes
B\'ezout's identity for the extended GCD underlying every modular inverse.
Runtime overflow is prevented by \texttt{\_ok} sentinel exports
(\texttt{ualbf\_compute\_sigma\_ok}, \texttt{ualbf\_mod\_inverse\_ok})
that verify results fit within 128 bits before the Rust engine reads the
truncated word halves.  The result is a system that executes
\emph{mathematically certified, C-linked binaries} inside a high-speed,
lock-free search engine.

\smallskip
In this paper, we present the following core contributions:

\begin{itemize}
  \item \textbf{Formalized Parity and Divisibility Lemmas:}
    We provide mechanically verified proofs in Lean~4 that a QPN cannot
    be a double square, and establish the foundational $\sigma$ parity
    and exact valuation lemmas (\cref{sec:math_foundations}).

  \item \textbf{Abundancy Index and Totient Ratio Analysis:}
    We formally verify that the abundancy index
    $H(N) = \sigma(N)/N$ is strictly bounded by the Euler totient ratio
    $N/\varphi(N)$, decompose this ratio into $H(N) \times C$ where $C$
    is a correction factor product, and prove $C < 1.022$ for QPNs
    coprime to~15, yielding $N/\varphi(N) < 2.0442$
    (\cref{subsec:abundancy_euler_ceiling,subsec:totient_decomposition}).

  \item \textbf{Cyclotomic Factorisation and Zsigmondy Obstructions:}
    We formalise the cyclotomic polynomial decomposition of
    $\sigma(p^{2e})$ and decompose the proof of Zsigmondy's theorem
    into seven modular sub-lemmas, establishing
    that high exponents spawn primitive prime divisors triggering the
    modulo-8 obstruction (\cref{subsec:cyclotomic_zsigmondy}).
    All sub-lemmas are now formally verified in Lean~4/Mathlib.
    The decomposition exposes every analytic ingredient
    as a separate lemma, creating a 100\% mechanically verified
    cyclotomic decomposition with zero gaps.

  \item \textbf{Standalone Correction Factor Module:}
    We provide a pure-$\mathbb{Q}$ telescoping-sum framework in
    \texttt{Pure/RationalBounds.lean} (the 36/35 pure identity machinery), 
    which serves as a direct mathematical dependency to bound the product 
    tail for the unified 2.0442 Euler ceiling theorem, requiring no real-analysis imports
    (\cref{subsec:correction_factor_standalone}).

  \item \textbf{Formal Exhaustion Certificate:}
    We prove the contrapositive
    \texttt{no\_\allowbreak{}solution\_\allowbreak{}no\_\allowbreak{}qpn}:
    if the Rust engine's ray-cast exhausts all candidate suffixes
    without finding $\sigma(N)=2N+1$, then $N$ is formally proven
    non-quasiperfect for the given prefix
    (\cref{subsec:bipartition}).

  \item \textbf{Theoretical Framework for Deeper Horizons: Exact Valuation Ray-Cast Strategy:}
    We introduce a ray-cast algorithm in Rust whose \emph{targeting
    step}---computing the unique residue class of $N_R$ modulo
    $\sigma(N_L)$ via a single modular inverse---runs in
    $\mathcal{O}(1)$ time. Candidates along the resulting arithmetic
    progression are then enumerated. However, we transparently note that for the
    $10^{\TelemetryMaxLog}$ bound, shallow structural traps completely exhausted the
    search space before this deep-horizon algorithm could activate
    (\cref{sec:computational_methodology}).

  \item \textbf{Extended Computational Bounds:}
    We successfully push the lower bounds for QPNs to $N > 10^{\TelemetryMaxLog}$,
    empirically validating our theoretical framework without
    encountering false positives (\cref{sec:results}).
\end{itemize}


%% -------------------------------------------------------------------
\subsection{The Quasiperfect Number Problem}
\label{sec:qpn_problem}

For a positive integer~$N$, the \emph{sum-of-divisors function} is
$\sigma(N) = \sum_{d \mid N} d$.  The \emph{abundancy index} of~$N$ is
the ratio $\sigma(N)/N$.  A \emph{perfect number} satisfies
$\sigma(N)=2N$; analogously, a \emph{quasiperfect number} (QPN) is a
positive integer satisfying
%
\begin{equation}\label{eq:qpn-def}
  \sigma(N) \;=\; 2N + 1.
\end{equation}

\begin{definition}[Quasiperfect Number]\label{def:qpn}
  A positive integer~$N$ is \textbf{quasiperfect} if $N > 0$ and
  $\sigma(N) = 2N + 1$.  In our Lean~4 formalisation this is encoded
  directly:
  \begin{center}
    \leaninline{def IsQuasiperfect (n : ℕ) : Prop := n > 0 ∧ sigma n = 2 * n + 1}.
  \end{center}
\end{definition}

No quasiperfect number has ever been found.  Cattaneo~\cite{cattaneo1951}
showed that any QPN must be an odd perfect square, and Hagis and
Cohen~\cite{hagis1982} established that $N > 10^{\TelemetryMinLog}$ with at least
seven distinct prime factors.  Because $2N+1$ is manifestly odd, one
obtains immediately that $\sigma(N)$ is odd for any QPN; the
consequences of this parity constraint are the starting point for the
analysis below.

\paragraph{Setup.}
The project setup explicitly relies on our foundational dependencies: Lean 4~\cite{moura2021}, Mathlib~\cite{mathlib2020}, Rayon~\cite{rayon2024}, bnum~\cite{bnum2024}, and ed25519-dalek~\cite{ed25519dalek2024}.
